
%% We're changing the margins for the longtable otherwise it'll be really long.
\newgeometry{left=1cm,right=1cm,top=1cm,bottom=2cm}
\begin{landscape}
\pagestyle{plain}
\footnotesize 
\renewcommand{\tabcolsep}{0.25cm}

% \begin{mdframed}[backgroundcolor=orangebg,
%                 rightline=false,
%                 leftline=false,
%                 topline=false,
%                 bottomline=false]
\begin{longtable}{P{0.7cm}
                  P{4cm}
                  P{6cm}
                  P{2cm}
                  P{4cm}
                  P{2cm}
                  P{3.5cm}}

\caption{UCC School of Computer Science and Information Technology Athena Swan Action Plan}\label{tab:actions}\\
\vspace{-57pt}
\hspace*{-9pt}
\crule{12mm}{3mm}\\
\toprule
    ID & 
    Action & 
    Rationale & 
    Responsibility &
    Outputs/Milestones &
    Time Frame & 
    Success measures\\
\midrule
\endfirsthead
\multicolumn{6}{l}%
{\textcolor{black}{%\tablename\ \thetable\ 
 %\textit{Continued from previous page}
}
} \\
\midrule
    ID & 
    {Action} & 
    {Rationale} & 
    {Responsibility} &
    {Outputs/Milestones} &
    {Time Frame} & 
    {Success measures}\\
\midrule
\endhead
%\midrule \\
%\multicolumn{6}{r}{\textcolor{black}{\textit{Continued on next page}}} \\
\endfoot
\bottomrule 
\endlastfoot



%%%%%%%%%%%%%%%%%%%%%%%%%%%%%%%%%%%%%%%%%%%%%%%%%%%%%%%%%%%%%%%%%%%%%%%%%%%%%
\midrule 
\multicolumn{7}{c}{\hyperlink{priority:1}{\textbf{Priority 1: Develop governance structures, policies  and procedures to ensure that the Athena Swan Ireland Principles are imbedded in the school}}}\hypertarget{p1}\\

\midrule 

\ref{action:intro:edi}\hypertarget{a2}
&
\multirow{3}{=}{\actintroedi} 
& 
A: To imbed the AS principles in the school's culture, governance, policies and procedures and to link with the wider EDI structures within the University. 
& 
Head of School
& 
Milestone 1: EDIIC committee established as part of the school governance structure.
& 
Mid 2023
&
Gender balanced EDIIC Committee established.\\

\cdashline{3-7}
\addlinespace 

&
& 
\multirow{2}{=}{B. The Chair of this committee will engage in ongoing consultation with the Head of School and the UCC EDI Unit to oversee the implementation of the AS action plan, and will sit on the College's AS steering committee.}
&
EDIIC Chair
&
Milestone 2: EDIIC Chair will monitor and oversee the implementation of the AS action plan.
&
2023, 2024, 2025, 2026
&
\multirow{2}{=}{Actions delivered based on time frame indicated; adjusted as needed.} \\

\cdashline{4-6}


& 
&
&
EDIIC Chair
&
Milestone 3: Annual report on implemented and outstanding actions.
&
2023, 2024, 2025, 2026\\




\midrule 


%%%%%%%%%%%%%%%%%%%%%%%%%%%%%%%%%%%%%%%%%%%%%%%%%%%%%%%%%%%%%%%%%%%%%%%%%%%%%



\ref{action:intro:agenda}\hypertarget{a5}
&
\actintroagenda
&
Ensures that the school is informed of progress with regard to implementation of the AS action plan, increase awareness of Athena Swan Ireland Principles and provide opportunity for feedback.
&
EDIIC Chair
& 
EDIIC Chair will report at school meetings.
& 
Twice yearly; October, April
&
The school reports satisfaction with advancement of EDI agenda based on the EDI staff survey in 2025.\\

\midrule 

\ref{action:intro:review:year3}\hypertarget{a6}
&
\actreviewimpactyearthree %% priority 1
&
To give staff and students an opportunity to give anonymous feedback on the impact of the AS action plan and to evaluate if the desired cultural changes are being brought about. To afford the EDIIC committee the opportunity to make any necessary mid-course correction.
&
EDIIC committee
&
Formal report to the school at the end of 2025.
&
November 2025
&
All key cultural metrics improved with respect to the original EDI survey.\\

\midrule

\pagebreak

\ref{action:record:hiring:phd}\hypertarget{aa0}
&
\multirow{3}{=}{\acttrackgenderinfophd}
&
\multirow{2}{=}{Research Centres and Principal Investigators currently manages these processes and data, but there is no central data source for PhD recruitment. Collating this data will provide a school-level view which will allow us to track trends across the school's PhD students such as differences by gender in completion rates/times.} 
&
PhD Programme Director
&
Milestone 1: A policy defined, approved by the school and implemented.
&
2023
&
\multirow{2}{=}{The school will have enhanced data as detailed in the annual PhD Programme Director's report that will inform strategies to improve gender balance.}\\

\cdashline{4-6}

&&&
PhD Programme Director
&
Milestone 2: Annual reports to the EDIIC/school detailing the gender-balance status in the PhD cohort.
&
December 2023, 2024, 2025, 2026 
&
\\

\cdashline{3-6}

&&
To monitor PhD applications and recruitment data to inform policies and procedures to improve gender balance across this cohort.  Currently the percentage of female PhD students range from  33\% in 2017-33\%, 19\% in 2018-28\% and 13\% in 2019. 
&
PhD Programme Director
&
Milestone 3: Database tracking gendered PhD progress and completions.
&
2023
&\\




\midrule

\ref{action:genderproof:guidelines}\hypertarget{b1}
&
\actgenderproof
&
The way the school presents itself influences how prospective staff and students from all genders and backgrounds evaluate the school as a place to work and study. As CS is a male dominated discipline, it tends to attract fewer female staff and students. Attracting female applicants at all levels remains a key challenge: Percentage female applicants to: 1) academic positions averages 19\% over past 3 years; 2) research posts, 20\%; 2) female UG registered, 24\%; 3) PhD students registered, 12\%.
&
Ad-hoc working group appointed by the Head of School
&
Milestone 1: Develop guidelines and training for staff on how to gender proof advertising, website, communications etc. for staff and student recruitment. Tools such as Gender Decoder, e.g. \url{https://capd.mit.edu/resources/gender-decoder/} may be used. Unconscious bias training amongst others will also be implemented.   
&
2023	
&
100\% of advertisements, all future materials and visuals follow the guidelines. \\

\cdashline{4-6}

&&&
Programme Directors & Milestone 2: all prospectus and web materials revamped to align with gender proof guidelines. 
& 2024 
&\\

\cdashline{4-6}
&&&
EDIIC	
&
Milestone 3: Annual verification compliance with gender-proofed guidelines. 
& December 2024, 2025, 2026 
&\\

\midrule

\ref{action:acad:cvgaps}\hypertarget{c6}
&
\actincludeleavestatement
&
During discussions in working and focus groups, a number of staff mentioned that there is no opportunity in job applications to include Covid-19 or family related leave statements to explain career gaps due to caring responsibilities. As this disproportionally affects women, applicants for academic, research, and professional and management posts should be able to provide relevant leave statements.
&
Head of School, School Manager
&
Guidelines will be created and communicated to the school on how provide allowance for relevant leave statements.
&
2023	
&
All advertisements will contain an invitation to candidates to provide a relevant leave statement. \\
	
\midrule 
\ref{action:acad:handbook}\hypertarget{cxx1}
&
\acthandbook %% priority 1
& 
From the focus group, research staff were unclear about opportunities to participate in tutoring, lecturing and other academic duties in the school.
&
Head of School and School Manager
&
Policy and application form for researcher participation in stated opportunities.
&
2024
&
10\% of researchers taking up one or more of the opportunities defined in the policy.\\

\midrule 
	


\ref{action:acad:transparentworkload}\hypertarget{c19}
&
\actreviewworkloadmodel
&
45\% of academic staff expressed dissatisfaction with the current workload model used in the school and belief it to be unfair. 
&
Head of School
&
The school workload model for academic staff will be reviewed to determine any source of dissatisfaction. 
&
2024; 2026	
&

Sources of dissatisfaction under the control of the school will be addressed, where possible. Staff satisfaction will be determined using the EDI annual survey in 2025, aiming to increase staff satisfaction by 25\%. \\ 

\midrule

\ref{action:pms:matchdescription}\hypertarget{d3}
&
\actreviewadvertisements
&
The research PMS templates currently used by the university for research recruitment are generic and used for both post-doctoral researchers and research support staff (PMS roles).  The templates require experience not relevant to a research support role, such as requesting publication records. This is likely to deter some candidates due to misinterpretation of the role. This drawback in the recruitment process has been highlighted by school PIs and corroborated by recent research PMS recruits. 
&
Head of School, School Manager
&
Guidelines for advertisements of research posts, and in particular research support posts, that provide practical advice on tailoring the duties and responsibilities associated with the research post.	
&
2023
&
Increased number of suitable applicants for PMS research posts, compared to 2020.\\


\midrule

\ref{action:cult:flexwork}\hypertarget{e5}
&
\actadoptflex
&
The introduction of blending working during the Covid-19 pandemic initiated the introduction of a UCC blended pilot working policy. Staff require a clear and transparent articulation of this policy and how it is implemented in the school.	
&
Head of School, line managers
&
A document outlining the implementation of university policies around blended working and flexible working hours at the school level, taking into consideration the requirements of different roles. 
&
Mid 2023, and updated as university policy evolves.
&
In the EDI survey 2025, staff report that the school implementation of the university policy is clear and transparent.\\


\midrule 




\pagebreak 


\multicolumn{7}{c}{\hyperlink{priority:2}{\textbf{Priority 2: Increase ratio of female applicants for staff posts (Academic/ Research) to improve staff gender  in these cohorts}}}\hypertarget{p2}\\
\midrule 


\ref{action:acad:recruittraining}\hypertarget{c7}
&
\actrecruitmenttraining
&
All staff engaged in recruitment for academic posts must have had recruitment training, however, this is not a requirement for research recruitment. Moreover, there has been no requirement for gender-balanced committees for research recruitment. In fact, the survey data suggests that these committees have not always been gender-balanced (19\% F, on average). Therefore, to align research recruitment and academic recruitment procedures, the school will introduce a new policy to cover staff training and to formalise guidelines for the operation of  research recruitment  committees.
&
Research Committee
&
A school policy on research recruitment, created and approved by the school. 

&
2024
&
All staff engaged in research recruitment have received appropriate training.

Research Recruitment Committees with be gender-balanced. Aspiring to 40\% female participation, where practicable.

\\


\midrule

\ref{action:acad:marketing}\hypertarget{c2}
&
\actimproveadvertisingtext
&
Attracting female applicants at all levels remains a key challenge as shown by the limited number of applicants. The percentage female applicants for academic positions averages 27\% over past 3 years.
&
Head of School, Centre Directors and PIs
&
Milestone 1: Gender-proofed advertising. See also Action~\ref{action:genderproof:guidelines}. \

Target a wide range of relevant channels to disseminate advertising material, such as Women in AI and IEEE Women in Engineering network.
&
2023
&
\multirow{3}{=}{Increase the percentage of female applicants for academic posts from an average of 27\% to 40\% by 2026. 
 Increase the percentage of female applicants research staff from 20\% to 40\% by 2026 }\\

\cdashline{4-6}

&&&
Head of School, School Manager. 	
&
Milestone 2: A record of the percentage of female applicants for academic positions.  	
&
2023, 2024, 2025, 2026.
&\\



\midrule 

\ref{action:acad:femaleonly}\hypertarget{c1}
&
\acttargetfemales
&
The survey data reveals that only 27\% of applicants for academic post, and 20\% for research posts,  have been female in the period from 2017 to 2020. We will actively target high-quality female applicants and will engage with targeted gender equality funding initiatives such as the SALI Professorships (See also UCC's Institutional Action~\hyperlink{act514}{5.1.4})
&
Head of School, Centre Directors and PIs.
&
A record of the percentage of female applicants for academic positions.  
&
2023, 2024, 2025, 2026
&
The number of female applicants for academic (27\%) and research (20\%) posts will be increased to 40\% by 2026, to align with university targets.
\\



\midrule 
\multicolumn{7}{c}{\hyperlink{priority:3}{\textbf{Priority 3: Increase number of female students registered in UG and PG programmes}}}\hypertarget{p3}\\
\midrule 

\ref{action:survey:females}\hypertarget{b0}
&
\actsurveyugfemales
&
A detailed understanding of perceptions of our courses and discipline among females is key to ensuring any interventions planned are rooted in reality. Attracting female applicants to UG programmes is a global challenge: Percentage females registered on UG programmes is - 24\%; but lower for our flagship BSc Computer Science - 14\% and BSc Data Science and Analytics - 20\%.  This will provide key insight in the lack of progression of female undergraduates to PG study. 	
&
Programme Directors	

&Survey data will provide information that can influnce our recruitment strategy with aim of attracting more females. 	
&
2023; 2025	
&
A female UG participation in the survey of 75\%\\



\midrule
\ref{action:review:outreach}\hypertarget{b2}
&
\actreviewactivities
&
Current activities aim to promote gender balance, by including females in web photos, videos, on open day stands, in brochures etc., however, there has been no follow up activity to determine the effectiveness of these measures.	
&
Ad-hoc working group, established by the Head of School.
&
The school have a set of guidelines for producing more effective female-friendly marketing materials and outreach activities.	
&
2023	
&
The percentage females entering BScCS and BScDSA will have increased from 14\% to 20\% and from 20\% to 25\%, respectively, by 2026.\\


\midrule

\ref{action:boost:visibility}\hypertarget{b3}
&
\actboostvisibility
&
Examples of successful females in the discipline will help to project computer science as a more female-friendly subject. 	
&
Outreach committee, seminar series coordinator
&
Seminars and talks and video presentations by high-profile female computing professionals. 
&
2023
&
25\% of seminar presenters will be female at school run seminars. \\

\midrule 

\ref{action:edi:coach}\hypertarget{b4}
&
\actappointmentor
&
Given the under representation of female undergraduate student (BSc Computer Science 14\%); BSc Data Science and Analytics (20\%), the school will advise and encourage female undergraduate students to avail of opportunities to support their career development and act as role models to attract more females into our programmes.	
&
Head of School
&
A programme of activities and supports for female undergraduates, such as, internships, workshops, and networking opportunities.
&
2023	
&
Increase  percentage of females entering BSc Computer Science to 20\% and BSc Data Science and Analytics to 25\% over the lifetime of the award.\\

\midrule 

\ref{action:enrich:opportunities}\hypertarget{b5}
&
\actenrichopportunities
&
In recent years, no female has progressed from CSIT undergraduate studies to postgraduate. In fact, in the past 5 years, only two students have done this, both male. By offering a number of female-only internships on an annual basis, it is hoped that more females will become familiar with research as a career option and will subsequently progress from undergraduate to postgraduate studies in CSIT.  
&
Head of School and Research Centre Directors
&
A female-only internship programme is established.

&
2024, 2025, 2026	
&
At least two female-internship opportunities are made available and taken-up on an annual basis. \\

\midrule

\pagebreak 

\multicolumn{7}{c}{\hyperlink{priority:4}{\textbf{Priority 4: Improved staff development to support for career progression}}}\hypertarget{p4}\\



\midrule

\ref{action:acad:mentorresearchers}\hypertarget{c4}
&
\actmentorresearchers
&
The pathway from postgraduate to career academic is challenging. A personal development plan can be a useful tool in helping aspiring academics to come to grips with the prerequisites of an academic post; and mentoring for the production of such a plan can mean the difference between success and failure. Given the underrepresentation of females in the discipline, it is felt that such mentoring could be extremely valuable to female researchers, since the percentage of female applicants to academic positions is far less than that of males. From the survey data, only 38\% of postdoctoral researchers currently have access to adequate mentoring.
&
Head of School, Research Centre Directors
&
A mentoring programme for creating a personal development plan is in place.
&
2023, 2024, 2025, 2026
&
80\% of postdoctoral researchers have a PDP, as measured by the EDI survey conducted in 2025.
				\\
				
				
\midrule

\ref{action:acad:mentoring}\hypertarget{c10}
&
\actmentorpromotion
&
From the survey data, only 33\% (50\% F, 17\% M) of academic staff reported access to adequate training and mentoring, to help them to navigate the promotion process. 
&
Head of School
&
A workshop for staff seeking promotion and a mentor appointed to interested staff following the PDRS review.
&
2023, 2024, 2025, 2026
&
100\% of staff are satisfied with mentoring scheme as measured by the EDI survey in 2025. \\
		
	



\midrule 

\ref{action:acad:pdrs}\hypertarget{c13}
&
\actrestartpdrs
&
The school has not operated a PDRS for over 5 years, although it was due to restart prior to the restrictions imposed by the pandemic in March 2020. In 2021, a new Head of School was appointed and restarting this process is a priority to support staff development.  	
&
Head of School, Line Managers.	
&
The PDRS in operation. 
&
2023, 2025
&
Relevant staff will have taken part in the PDRS by the end of 2023. \\

\midrule 

\ref{action:acad:supportfunding}\hypertarget{c16}
&
\actsupportprogramme
&
The school has an excellent track record in attracting research funding, however, this is centered around a small number of senior staff. The survey revealed several issues in relation to the opportunities to apply for funding, the support from the school in the process, but also for those whose applications were unsuccessful, and support from university.
&
Research Committee
&
Workshop to support interested staff navigate the funding process and a peer-review process to appraise applications before submission.
&
2023, 2024, 2025, 2026.
&
100\% of applicants are satisfied with the supports provide by the school for research funding applications,  as measured by the EDI  staff survey in 2025. \\




\midrule 
\ref{action:pms:experience}\hypertarget{d8}
&
\actpmsopportunities
&
School PMS staff expressed a desire, in the EDI survey and focus group,  to develop skills beyond the remit of their existing roles. A range of measures were discussed including, committee membership, training, deputising where appropriate and the taking on of additional responsibilities.	
&
PMSS Line Managers	
&
A record of occurrences of opportunities offered.
&
2023, 2024, 2025, 2026
&
A minimum of 2 opportunities are offered annually and school PMS staff record satisfaction with same as part of the EDI survey in 2025.\\

\midrule


\multicolumn{7}{c}{\hyperlink{priority:5}{\textbf{Priority 5: Promote a culture of inclusivity and belonging}}}\hypertarget{p5}\\
\midrule 


\ref{action:cult:diverseoutreach}\hypertarget{b6}
&
\actreviseoutreachpractice
&
Discussions in focus groups highlighted the need to increase diversity in the secondary school outreach programme. Members felt strongly that increasing diversity in education promotes a culture of inclusion and leads to better outcomes for students from minority-backgrounds.	
&
Outreach Committee	
&
A list of schools identified as having students with diverse backgrounds. 
&
2023, 2024, 2025, 2026 
&
At least 20\% of secondary school visited annually were to those schools identified as having students with diverse backgrounds by 2026.	\\







\midrule 

\ref{action:cult:researchinteraction}\hypertarget{e2}
&
\actimproveinteraction
&

During focus group sessions researchers and PhD students expressed a feeling of separation from the school. They believed that a lack of school-wide events reduced opportunities for collaboration and inclusion.
&
Research Committee and Director of PhD Programme.
&
School research forum organised annually
Social event organised after the annualy PhD review day for PhD students and staff. 
Invitations issued to reserachers to attend the monthly school coffee morning. 
On arrival, new research staff and PhD students are introduced to the staff via the weekly ``Monday Morning News'' email.
&
2023, 2024, 2025, 2026
&
75\% of researcher staff and PhD students report feeling included in the EDI staff survey in 2025.\\










\midrule
\ref{action:cult:backgrounds}\hypertarget{d10}
&
\multirow{3}{=}{\actcultbackground} %% priority 
& 
\multirow{3}{=}{We are not aware of the needs of our students from diverse backgrounds nor have we systematically engaged with issues of equality and diversity to enhance student experience within the curriculum.}
& 
TLSE Committee
&
Milestone 1: Determine needs of students with diverse backgrounds using the  biannual EDI Student Survey
&
2023
&
\multirow{3}{=}{10\% of academic staff participating in the EDI/IHREC briefings and training; and 5\% of academic staff engaging with the DISC toolset in 2024.}\\
\cdashline{4-6}

&&&
TLSE Committee
&
Milestone 2: Explore and exploit appropriate tools and training to enhance the student experience.
& 2024 
&\\

\cdashline{4-6}
&&&
TLSE Committee
&
Milestone 3: Participate in the IHREC eLearning module on ‘Equality and Human Rights in the Public Service’ and the UCC’s EDI in HE programme
& 2024 
&\\


\midrule

\ref{action:cult:famleave}\hypertarget{e5}
&
\actremoveimpediments
&
Feedback from the staff survey and focus groups reported a lack of clarity on how family leave is managed in the school. Researchers described challenges requesting family leave and principal investigators voiced concern over different approaches to funding family leave by different research funders.
&
Head of School, Line Manager
&
A policy outlining how family leave arrangements are managed by the school within the constraints of statutory obligations.
&
2024
&
Staff report clarity and understanding on how the school implements family leave in the EDI survey of 2025.

\\


\end{longtable}
% \end{mdframed}


\end{landscape}
\restoregeometry